% !TEX root = main.tex

% ============================================================================
% FRONT MATTER
% ============================================================================
% Roman numeral page numbering for all front pages.
% Title page is counted as page i but number is suppressed.
% Arabic numbering resumes at the end of this file.

\pagestyle{empty}
\pagenumbering{roman}

% ============================================================================
% T I T L E   P A G E
% ============================================================================

\begin{titlepage}

\newgeometry{left=0.8in, right=0.8in, top=1in, bottom=1in}
\centering

{\large \univname}

\vfill

{\Huge \bfseries Dynamic RAG under Energy Constraints: \\
Characterizing the Structural\\ \vspace{0.2cm}
Limits of Module-Level Routing}

{\large\itshape \tcontext}

\vspace{0.5cm}

{\normalsize by}

\vspace{0.5cm}

{\large \href{mailto:vieirrui@students.zhaw.ch}{Rui Vieira}}

\vfill

{\normalsize Advisor: \supnameA\\Supervisor: \supnameB}

\vspace{0.5cm}

{\normalsize Host Institution:\\
LLM Lab -- Center for IoT Innovation, \\
National Taiwan University of Science and Technology}

\vfill

\begin{minipage}{0.6\textwidth}
  \centering\small
  Thesis submitted in partial fulfillment of the requirements\\
  for the degree of Master of Science in Engineering
\end{minipage}

\vspace{0.5cm}

{\normalsize
  8 March 2026
}

\restoregeometry
\end{titlepage}

% Front pages: plain style with Roman numerals starting at ii
\pagestyle{plain}
\setcounter{page}{2}

% ============================================================================
% E X A M I N I N G   C O M M I T T E E
% ============================================================================

\cleardoublepage
\begin{center}\textbf{\Large Examining Committee Membership}\end{center}
\bigskip

\noindent
The following served on the Examining Committee for this thesis.

\bigskip

\noindent
\begin{tabbing}
Supervisor:\quad \= \kill
Advisor: \> \supnameA \\
\> \supinstituteA \\
\> School of Engineering, Zurich University of Applied Sciences \\
\> \href{mailto:\supmailA}{\supmailA}
\end{tabbing}

\bigskip

\noindent
\begin{tabbing}
Supervisor:\quad \= \kill
Supervisor: \> \supnameB \\
\> Center for IoT Innovation \\
\> Department of Industrial Management \\
\> National Taiwan University of Science and Technology \\
\> \href{mailto:\supmailB}{\supmailB}
\end{tabbing}

% ============================================================================
% A B S T R A C T
% ============================================================================

\cleardoublepage
\begin{center}\textbf{\Large Abstract}\end{center}
\bigskip

Chatbot assistants powered by large language models (LLMs) have seen
rapid adoption, yet the standard Retrieval-Augmented Generation (RAG)
architecture relies on ``always-on'' modular pipelines that are computationally
inefficient. This thesis investigates RAG pipelines from a system perspective,
treating energy consumption as a first-class metric to determine if selective
module activation can reduce operational costs without compromising answer
quality.

To this end, an energy measurement framework was developed for a distributed
four-node RAG system deployed as a knowledge management platform at the Center
for IoT Innovation (CITI), National Taiwan University of Science and Technology.
The framework attributes request-level energy, latency, and quality across
infrastructure nodes and pipeline stages. Empirical analysis of 1,200 multi-hop
queries from the HotpotQA benchmark reveals that always-on configurations are
significantly inefficient: the full pipeline consumes 5.6$\times$ more energy
than a minimal baseline for a marginal quality improvement of less than 0.06
points in factual correctness. Notably, the HyDE (Hypothetical Document Embeddings)
module accounts for the largest energy overhead while actually degrading quality in
the studied dense-retrieval environment.

To exploit this cost-quality asymmetry, we formalize per-query configuration control
as an optimization problem. While an oracle router could achieve a 75\% energy savings
ceiling, a systematic evaluation of five routing strategies (including multiclass
classification and utility scoring) reveals a structural limitation. All pre-execution
routing models based on query-only features converge to a configuration accuracy of
39.5\%, failing to meaningfully outperform a static ``cheapest-acceptable'' baseline.
This suggests a circular dependency in RAG optimization: the information required
to predict module benefit often emerges only after retrieval and intermediate
execution; the very computation the router seeks to skip. These findings shift the
focus for RAG efficiency from complex predictive modeling to policy-driven defaults
and lightweight post-retrieval probing.

\vspace{1cm}
\textbf{Keywords:} retrieval-augmented generation, energy efficiency, query routing,
LLM inference, conditional computation, local deployment, RAG optimization.

% ============================================================================
% A C K N O W L E D G M E N T S
% ============================================================================

\cleardoublepage
\begin{center}\textbf{\Large Acknowledgments}\end{center}
\bigskip

% TODO: Write acknowledgments
This is the only page throughout this thesis where I am not required to be a neutral observer, 
so I will take the opportunity to be sincere, if brief.

First, I must express my deepest gratitude to my advisor at ZHAW, Dr. Stefan Czerner. 
Thank you for your unwavering support and for the constant interest you showed in this research. 
I am especially grateful for the flexibility to pursue a path that didn't strictly suit our local 
lab development, even when that meant supporting my plan to ``run away'' to the other side of 
the world. Your positive attitude made the challenges feel manageable and the advice you shared 
is something I will try to carry with me for as long as possible.

I also want to thank the ZHAW School of Engineering for the institutional support and for granting 
me the unique opportunity to conduct this research abroad.

My time in Taiwan would not have been possible without the vision of Distinguished Prof. Shuo-Yan Chou. 
Thank you for welcoming me into the Center for IoT Innovation (CITI) and for your generosity in fostering 
such a vibrant, diverse and skillful research environment. Your ability to identify what truly makes research 
impactful has been a lesson in itself.

To the "lab rats" at CITI, especially Harry, William and Aesop for constantly helping out and the 
S-TIER LLM Team (Avi, Ishaq, Zhafran and Sir Fadhil): thank you for the long days and the even 
longer nights. You were my teachers, my IT support and my back. Thank you for the laughter, the 
hard-to-describe humor and the (already forgotten) Mandarin and Indonesian swear words. Knowing 
I could show up every day to your smiles made all the difference.

Obviously, to my friends and my family, who are the kindest and prettiest souls I know: you made me who I am. 
Thank you for believing in me, often more than I believe in myself.

Aos meus pais e à minha irmã, que me apoiaram sempre. Embora eu não diga o suficiente, saibam que todo o 
meu esforço é também por vocês. Provavelmente vocês não lerão esta tese, mas se por acaso virem estas palavras, 
quero que saibam que amo-vos muito.

% ============================================================================
% T A B L E   O F   C O N T E N T S
% ============================================================================

\cleardoublepage
\phantomsection
\setcounter{tocdepth}{2}
\renewcommand\contentsname{Table of Contents}
\tableofcontents

% ============================================================================
% L I S T   O F   A B B R E V I A T I O N S
% ============================================================================

\cleardoublepage
\section*{List of Abbreviations}

\noindent
\begin{tabular}{@{} l @{\hspace{2em}} l @{}}
API     & Application Programming Interface \\
AUC     & Area Under the Curve \\
BGE     & BAAI General Embedding \\
BM25    & Best Matching 25 \\
CITI    & Center for IoT Innovation \\
CPU     & Central Processing Unit \\
EM      & Exact Match \\
GPU     & Graphics Processing Unit \\
HyDE    & Hypothetical Document Embeddings \\
ICC     & Intraclass Correlation Coefficient \\
IoT     & Internet of Things \\
J/tok   & Joules per Token \\
KILT    & Knowledge-Intensive Language Tasks \\
KMS     & Knowledge Management System \\
LLM     & Large Language Model \\
LoRA    & Low-Rank Adaptation \\
LR      & Logistic Regression \\
MLP     & Multilayer Perceptron \\
NLP     & Natural Language Processing \\
NVML    & NVIDIA Management Library \\
RAGAS   & RAG Assessment \\
RAG     & Retrieval-Augmented Generation \\
RAPL    & Running Average Power Limit \\
RF      & Random Forest \\
RRF     & Reciprocal Rank Fusion \\
TTFT    & Time to First Token \\
vLLM    & Virtual Large Language Model inference engine \\
XGBoost & Extreme Gradient Boosting \\
\end{tabular}

% ============================================================================
% L I S T   O F   F I G U R E S
% ============================================================================

\cleardoublepage
\listoffigures

% ============================================================================
% L I S T   O F   T A B L E S
% ============================================================================

\cleardoublepage
\listoftables

% ============================================================================
% Switch to Arabic numbering for main content
% ============================================================================

\cleardoublepage
\pagenumbering{arabic}
